\section{Introduction}

A \emph{standard flip} \(X\dashrightarrow X'\), in the setting of
Bondal--Orlov and of Belmans, Fu and Raedschelders \cite{BFR}, is built from
the following data.  There is a birational morphism \(p\colon X\to\bar X\)
of smooth projective varieties, an isomorphism away from \(F=p^{-1}(Z)\) for
a smooth \(Z\subset\bar X\), and vector bundles \(V,V'\to Z\) of ranks \(r\)
and \(s\) with
\[
  F\cong\PP(V),\qquad
  N_{F|X}\cong\pi^{*}(V')\otimes\mathcal O_{\PP(V)}(-1),
\]
where \(\pi\colon F\to Z\).  Blowing up \(F\) and contracting the exceptional
divisor \(E\cong\PP(V)\times_{Z}\PP(V')\) in the other direction produces
\(X'\), which contains \(F'\cong\PP(V')\) of codimension \(r\).  Throughout,
\(r>s\).  Since \(F\subset X\) has codimension \(s\) and \(F'\subset X'\) has
codimension \(r\), the two blow-up formulas for the canonical class on
\(\widehat X=\mathrm{Bl}_{F}X\) give
\(\tau^{*}K_{X}-\tau'^{*}K_{X'}=(r-s)E\).  We call
\[
  \nu:=r-s
\]
the \emph{discrepancy} of the flip; \(s=1\) is the blow-up of \(\bar X\)
along \(Z\), and \(s=0\) is the degenerate case \(X=F=\PP(V)\),
\(X'=\varnothing\), of a projective bundle.

Shen and Shoemaker \cite{SS} compute the spectrum of quantum multiplication
by \(c_{1}(X)\) after restricting the Novikov parameter to the extremal ray
contracted by \(p\), and prove that the Gamma-class decomposition of
\(H^{*}(X)\) induced by the semiorthogonal decomposition of \(\Db(X)\) of
\cite{BFR} is a decomposition into strong asymptotic classes in the sense of
Galkin--Golyshev--Iritani \cite{GGI}.  This is
\cite[Theorems~1.2, 1.4, 9.9 and~9.14, Corollary~1.5]{SS}; apart from
\cite[Theorem~1.2]{SS}, which assumes \(r-s>1\), it is stated for
standard flips with the explicit remark that it covers the blow-up case
\(s=1\) and the projective-bundle case \(s=0\).

The proof passes through an explicit hypergeometric series.  In the local
model \(T=\mathrm{tot}(\pi^{*}V'\otimes\mathcal O_{\PP(V)}(-1))\), after the
extremal specialization \cite[(33)]{SS} of Novikov and parameter variables,
\cite[Theorem~4.4]{SS} states that the extremal \(J\)-function of \(T\) is
\begin{equation}
  \label{eq:ss35}
  J^{T}(q,t,z)
  = z\,e^{t/z}q^{c_{1}(T)/z}
    \sum_{d\ge0}q^{(r-s)d}\,
    \frac{\prod_{j=1}^{s}\prod_{m=0}^{d-1}(\sigma_{j}-H-mz)}
         {\prod_{i=1}^{r}\prod_{m=1}^{d}(\rho_{i}+H+mz)},
\end{equation}
which is \cite[(35)]{SS}.  Everything after Section~4 of \cite{SS} runs on
this formula: the ring presentation \cite[(37)]{SS} is deduced from it, the
extremal spectrum is read off from that presentation, and the modified
\(J\)-function of \cite[(39)--(40)]{SS} used in the asymptotic analysis of
Sections~7--9 is defined from it.

Three hypotheses in the printed argument constrain \(\nu\).
\cite[Theorem~4.4]{SS} is stated under \(r-s>1\).
\cite[Remark~4.5(3)]{SS} asserts that for \(r-s\le1\) the right-hand side of
\eqref{eq:ss35} is not of the form \(z\one+(\text{lower order in }z)\) and
therefore cannot equal the \(J\)-function of \(T\), although it does lie on
Givental's Lagrangian cone of \(T\) and so is an \(I\)-function.  Finally,
the Barnes asymptotic expansion of the Meijer \(G\)-function of
\cite[Lemma~7.4]{SS} is introduced by the clause ``For \(r-s>1\)'', and the
sector \cite[(64)]{SS} printed with it is derived under that inequality.
Under the standing assumption \(r>s\), the range excluded by these three
hypotheses is exactly \(\nu=1\), that is
\[
  r=s+1 .
\]
It contains every codimension-two blow-up, the case \((r,s)=(2,1)\).  Thus
\cite[Theorem~1.2]{SS} is proved only in its stated range \(r-s>1\), while
\cite[Theorems~1.4 and~9.14]{SS} and \cite[Corollary~1.5]{SS}, which are
stated for standard flips and blow-ups without a restriction on \(\nu\), rely
on inputs whose printed proofs do not cover \(r-s=1\).

This note supplies the two missing steps.  The second uses only
\cite[Theorems~A.1 and~A.2]{SS}, in the form the authors state them.  The first uses
\eqref{eq:ss35} together with the fact that it lies on Givental's Lagrangian
cone.  That fact is asserted in \cite[Remark~4.5(3)]{SS}, the same remark
whose other half we contradict, so we do not quote it: Section~3 derives it
from the sources \cite{Brown} and \cite{CoatesGivental} that the authors
name, and checks that none of them restricts \(r-s\).

\begin{theorem}
  \label{thm:normalization}
  Let \(r=s+1\) with \(s\ge1\).  Then, in the extremal specialization
  \cite[(33)]{SS} and with the same formal variables and cohomological
  completion as in \cite[Theorem~4.4]{SS}, the degree-\(d\) summand of
  \eqref{eq:ss35} has \(z\)-order at most \(-1\) for every \(d\ge1\).
  Consequently \eqref{eq:ss35} equals \(z\one+\bar{\mathbf t}+O(z^{-1})\)
  with \(\bar{\mathbf t}=t+\ln(q)c_{1}(T)\), and equals the extremal
  \(J\)-function of \(T\).  In particular the presentation \cite[(37)]{SS}
  of \cite[Theorem~4.6]{SS} holds for \(r=s+1\),
  \(s\ge1\), as does \cite[Theorem~1.2]{SS}, with the single nonzero
  eigenvalue \(\lambda_{0}(q)=(-1)^{s}q\) of multiplicity
  \(\rk H^{*}(Z)\).
\end{theorem}

The negative assertion of \cite[Remark~4.5(3)]{SS} is therefore accurate, in
the range \(r>s\) it addresses, only for \(s=0\), where \((r,s)=(1,0)\) and
the degree-one summand of \eqref{eq:ss35} does contribute in \(z\)-order
zero.

\begin{theorem}
  \label{thm:aperture}
  Let \(r=s+1\) with \(s\ge1\), and let \(k\in\{0,1\}\).  The asymptotic
  expansion of \cite[Proposition~7.5]{SS} for the component indexed by
  \(m=-k\) is valid on
  \(\bigl|\arg(z/q)+(2m+s)\pi\bigr|<\tfrac{3\pi}{2}\); \cite[Theorem~A.1]{SS}
  does not supply it on the wider sector obtained by setting \(r-s=1\)
  in \cite[(64)]{SS}.  The expansion of
  \cite[Proposition~8.2]{SS} is valid on
  \(\bigl|\arg(z/q)-(1-s)\pi\bigr|<\tfrac{3\pi}{2}\).  These two sectors meet
  in an open sector of opening \(2\pi\) containing both the eigenvalue ray
  \(\arg(z/q)=-(2m+s)\pi\) and the tame ray \(\arg(z/q)=(1-s)\pi\).
\end{theorem}

For a codimension-two blow-up, \((r,s)=(2,1)\) and \(k=m=0\), the two
sectors are \(-\tfrac{5\pi}{2}<\arg(z/q)<\tfrac{\pi}{2}\) and
\(|\arg(z/q)|<\tfrac{3\pi}{2}\), and they meet in
\(-\tfrac{3\pi}{2}<\arg(z/q)<\tfrac{\pi}{2}\), which contains the tame ray
\(0\) and the eigenvalue ray \(-\pi\).

\begin{corollary}
  \label{cor:repaired}
  \cite[Theorems~1.2, 1.4, 9.9 and~9.14]{SS} and \cite[Corollary~1.5]{SS}
  hold for every standard flip with \(r=s+1\) and \(s\ge1\), in particular
  for every codimension-two blow-up.
\end{corollary}

The mechanism behind Theorem~\ref{thm:normalization} is a single count, and
it is worth saying where the natural first count goes wrong.  The numerator
of the degree-\(d\) summand of \eqref{eq:ss35} has \(sd\) factors and the
denominator has \(rd\) factors.  Treating all of them as linear in \(z\)
gives \(1+sd-rd=1-\nu d\) for the maximal \(z\)-order, which is \(0\) at
\(d=1\) when \(\nu=1\): a curve-induced contribution in \(z\)-order zero,
hence a nontrivial mirror map, hence exactly the failure asserted in
\cite[Remark~4.5(3)]{SS}.  But the inner index in the numerator runs over
\(0\le m\le d-1\), so \(s\) of its factors are \(\sigma_{j}-H\), free of
\(z\).  The correct maximal order is \(1+s(d-1)-rd=1-s-\nu d\), and at
\(\nu=1\) the \(s\) missing powers of \(z\) are precisely what pushes every
\(d\ge1\) summand below \(z^{0}\) as soon as \(s\ge1\).  Once no
\(z\)-order-zero contribution survives, the \(z^{0}\) coefficient of
\eqref{eq:ss35} is the parameter \(\bar{\mathbf t}\) itself, the mirror map
is the identity, and uniqueness of the \(J\)-slice of the cone converts cone
membership into an identification.

Theorem~\ref{thm:aperture} is not a strengthening but a correction of an
aperture.  The Meijer \(G\)-function of \cite[Lemma~7.4]{SS} has
\(\nu=r-s\), and \cite[Theorem~A.1]{SS} --- Barnes' expansion in the form
of Meijer \cite{Meijer} and Luke \cite{Luke} --- is valid on
\(|\arg t|<(\nu+\epsilon)\pi\) with \(\epsilon=1\) for \(\nu>1\) but only
\(\epsilon=\tfrac12\) for \(\nu=1\).  The sector \cite[(64)]{SS} is the
\(\epsilon=1\) sector; substituting \(r-s=1\) into it claims an aperture
\(\pi/2\) wider on each side than Barnes' expansion supports.  The point of
Theorem~\ref{thm:aperture} is that the smaller, correct sector is still
large enough: the common sector with \cite[Proposition~8.2]{SS} retains
opening \(2\pi\), so the oriented comparison of the nonzero-eigenvalue and
tame expansions that the argument of \cite[Section~9]{SS} performs is
available at \(\nu=1\) as it is for \(\nu>1\).

What is not covered is the formal endpoint \((r,s)=(1,0)\), the rank-one
projective bundle, where \(\PP(V)=Z\), \(X'=\varnothing\), and the point
fibres contain no line class \([L]\).  Its displayed hypergeometric expression
also fails the required order count.  Nothing here bears on \(\nu>1\), where
the argument of \cite{SS} already applies, and nothing here changes any
statement of \cite{SS} outside the two hypotheses named above.

Section~\ref{sec:hypotheses} fixes notation and records the three places
where \(\nu\) is constrained in \cite{SS}, together with what depends on
each.  Section~\ref{sec:normalization} proves
Theorem~\ref{thm:normalization}, Section~\ref{sec:aperture} proves
Theorem~\ref{thm:aperture}, and Section~\ref{sec:scope} assembles
Corollary~\ref{cor:repaired} and states the boundary of the repair.

\medskip
\noindent\emph{Convention.}  Numbered items of \cite{SS} are always cited
explicitly, as in \cite[Theorem~4.4]{SS} or \cite[(35)]{SS}, and refer to
the version arXiv:2502.08762v2.  Displays numbered here, such as
\eqref{eq:ss35}, are our own.
